% Modele : title-qr - page de titre, photo + QR (image) centres sous le nom
% Champs optionnels : title_tex, authors_formatted, event_tex, date_tex, location_tex, photo_path, qr_image_path
%
% Marge de securite : 1.4cm a gauche/droite (comme le bloc infos evenement).
% Marge du haut : 0.6cm, alignee sur la convention deja en place dans les autres
% gabarits du deck (plain/tableau/equation/remerciements). Tous les
% blocs sont en flux normal (pas de tikz overlay absolu) : la position verticale de
% chaque bloc depend uniquement de la hauteur cumulee des blocs precedents, ce qui
% evite tout risque de chevauchement entre un bloc en flux et un bloc ancre en absolu
% (bug observe et corrige a la tache 11 - le bloc auteur chevauchait le bloc infos
% evenement quand ce dernier etait un tikz overlay independant du flux).
%
% Deux dimensions distinctes (tache 12, correction de la tache 11) :
%   - ordre d'apparition vertical : bloc 1 (titre) > bloc 3 (nom+titre auteur)
%     > bloc 2 (photo+QR) > bloc 4 (infos evenement) - INCHANGE depuis la version
%     d'origine, ne pas reordonner selon l'importance.
%   - ordre d'importance communicationnelle : bloc 1 > bloc 2 > bloc 3 > bloc 4 -
%     sert uniquement a doser l'emphase/l'espacement, pas la position verticale.
% \NouBlocGap : meme valeur au-dessus et en-dessous du bloc 2 (photo+QR) dans
% l'ordre d'apparition ci-dessus, donc entre bloc 3->bloc 2 et bloc 2->bloc 4.
\newcommand{\NouBlocGap}{1.0em}

\begin{frame}[plain]
<%- background_block %>

\vspace{0.6cm}
\hspace*{1.4cm}\begin{minipage}{\dimexpr\paperwidth-2.8cm\relax}
  {\NouTitleFont\fontsize{21}{21}\selectfont\color{NoumanityBlack}\raggedright <%- title_tex %>\par}
\end{minipage}

\vspace{\NouBlocGap}
\hspace*{1.4cm}\begin{minipage}{\dimexpr\paperwidth-2.8cm\relax}
  {\small\color{NoumanityBlack}\raggedright <%- authors_formatted %>\par}
\end{minipage}

\vspace{\NouBlocGap}
\begin{center}
<% if photo_path and photo_path ~= "" then %>\includegraphics[height=1.9cm]{<%- photo_path %>}<% end %><% if photo_path and photo_path ~= "" and qr_image_path and qr_image_path ~= "" then %>\hspace{1.5em}<% end %><% if qr_image_path and qr_image_path ~= "" then %>\includegraphics[height=1.9cm]{<%- qr_image_path %>}<% end %>
\end{center}

% yshift=0.5cm de l'ancien overlay -> ~1.2em de vspace en flux : ordre de grandeur
% similaire, valeur minimale par defaut pour la marge du bas (tache 10 : mesuree a
% ~0.2cm avant correctif, jugee trop petite). \scriptsize compacte le bloc.
\vspace{0.6em}
\hspace*{1.4cm}\begin{minipage}{0.85\paperwidth}
  \raggedright
  {\scriptsize\color{NoumanityBlack!60}
    <%- event_tex %> \textbullet{} <%- date_tex %>\par
    \vspace{0.05em}
    <%- location_tex %>}
\end{minipage}

<% if brand_logo_path and brand_logo_path ~= "" then %>
\begin{tikzpicture}[remember picture, overlay]
  \node[anchor=south east, inner sep=0pt, xshift=-1.4cm, yshift=0.5cm]
    at (current page.south east) {%
      \begin{minipage}{5cm}
        \raggedleft
        \includegraphics[height=1.0cm]{<%- brand_logo_path %>}\par
        <% if brand_tagline_tex and brand_tagline_tex ~= "" then %>\vspace{0.15em}{\scriptsize\color{NoumanityBlack!70} <%- brand_tagline_tex %>}<% end %>
      \end{minipage}};
\end{tikzpicture}
<% end %>
\end{frame}
